\begin{frame}{이번 주는 어떻게 진행되는가}
  \small
  \begin{tabular}{@{}clp{0.4\textwidth}@{}}
    \toprule
    순서 & 내용 & 시간 배분 \\
    \midrule
    1 & 다른 두 팀의 발표 청취(팀당 5$\sim$7분, 16.1 형식) + 보고서·
      코드시트 열람 & 약 30분 \\
    2 & \textbf{작성 리뷰 2통} — 팀당 체크리스트 5항목 + 반박 가능한
      질문 $\geq 1$개(개념의 장·절 명시) & 발표 직후(각 10분) \\
    3 & 질의·토론 — 리뷰어가 쓴 질문에 발표팀이 그 자리에서 답변 &
      나머지 시간 \\
    \bottomrule
  \end{tabular}
  \vskip0.8em
  \begin{block}{핵심 = 2번의 \textbf{작성 리뷰}}
    3번 토론은 ``발표팀이 그 질문에 \emph{실제로} 답할 수 있는가''를 가르는
    자리 — \textbf{새로운 계산·실험}(시드 재실행, 평가 프로토콜 변경,
    베이스라인 실측, 하이퍼파라미터 민감성 실험)을 꺼내야 답하는 질문이
    \textbf{3점}이고, ``네, 알겠습니다''로 넘겨질 질문은 2점에 머문다.
    16.1에서 ``수식적 정당화''가 발표팀의 \emph{의무}였다면, 이번 주
    그 \emph{타당성}을 리뷰어가 \emph{심사}하는 것.
  \end{block}
\end{frame}
